% Generated draft: requires mathematical and source review.
\documentclass[11pt]{article}
\usepackage[margin=25mm]{geometry}
\usepackage{fontspec}
\setmainfont{Latin Modern Roman}
\usepackage{amsmath,amssymb,mathtools,unicode-math}
\setmathfont{Latin Modern Math}
\usepackage{enumitem,needspace,xurl,hyperref}
\setlength{\parindent}{0pt}
\setlength{\parskip}{5pt}
\setlength{\emergencystretch}{3em}
\newcommand{\N}{\mathbb{N}}
\newcommand{\Z}{\mathbb{Z}}
\newcommand{\Q}{\mathbb{Q}}
\newcommand{\R}{\mathbb{R}}
\newcommand{\C}{\mathbb{C}}
\newcommand{\F}{\mathbb{F}}
\newcommand{\A}{\mathbb{A}}
\newcommand{\PP}{\mathbb{P}}
\newcommand{\E}{\mathbb{E}}
\newcommand{\eps}{\varepsilon}
\newcommand{\dd}{\,\mathrm{d}}
\begin{document}
\section*{UnsolvedMath ID: 30004342}
\textbf{Source problem code:} catalogue entry
\subsection{Definitions and mathematical statement}
Let \(k,j\in\Z\) satisfy
\[
 k\geq 3,\qquad j\geq 0,\qquad j\equiv 2\pmod 4.
\]
The supplied record uses \(S_{2k+j-2}(\Gamma_1)\) for a space of cusp forms of weight \(2k+j-2\), but does not define \(\Gamma_1\). For a Hecke eigenform \(f\) in this space, it denotes by \(\phi_0\) a corresponding degree-one Jacobi form, and by \(E(\phi_0)\) its degree-two Klingen lift. The correspondence, Jacobi-form spaces, weights, indices, levels, and normalizations are not specified.

\textbf{Reported conjectural assertion.} For every such \(f\), there exist a degree-two holomorphic Jacobi form \(\Phi\) outside the lift and a prime ideal \(\mathfrak l\) such that
\[
 \mathfrak l\mid L^{\mathrm{alg}}(k+j,f)
 \quad\text{and}\quad
 \lambda_{\Phi}(T)\equiv
 \lambda_{E(\phi_0)}(T)\pmod{\mathfrak l}
 \qquad\text{for every integral Hecke operator }T.
\]
Here \(\lambda_G(T)\) denotes the eigenvalue in \(TG=\lambda_G(T)G\), whenever \(G\) is an eigenform for \(T\). Thus the assertion requires \(\Phi\) and \(E(\phi_0)\) to possess the stated simultaneous Hecke eigenvalues. This display records the quantifier order of the supplied statement: the prime ideal is existentially quantified together with \(\Phi\).

\textbf{Unresolved qualifications.} This is a formulation template requiring additional source data. The record does not specify the algebraic normalization of \(L^{\mathrm{alg}}(k+j,f)\), a common coefficient field, the ring containing \(\mathfrak l\), integrality conditions, or the integral Hecke algebra. If a common integral coefficient ring \(A\) is specified, the displayed congruence means that the eigenvalue difference belongs to \(\mathfrak l\subset A\) for each allowed \(T\); the record does not provide such a ring. The phrase ``outside the lift'' does not identify an excluded image or subspace. It is also unclear whether the intended assertion assumes in advance that a suitable prime divides the algebraic special value, or concerns every prime satisfying further arithmetic restrictions. These alternatives cannot be substituted for the supplied existential formulation without further evidence.

The supplied curation reports a proof under explicit arithmetic hypotheses in the report, but does not state those hypotheses. Consequently neither that conditional theorem nor a complete unconditional conjecture can be recovered precisely from this record alone.
\subsection{Short English statement}
For integers \(k\geq3\) and \(j\geq0\) with \(j\equiv2\pmod4\), the supplied assertion associates to each indicated cusp eigenform a degree-two holomorphic Jacobi eigenform outside an unspecified lift, congruent to its Klingen lift modulo a prime dividing an algebraic special \(L\)-value. Essential spaces, normalizations, prime restrictions, and arithmetic hypotheses remain unspecified. The supplied partially solved label is dated curation, not an independently verified conclusion.
\subsection{Sources}
\begin{itemize}
\item Local UnsolvedMath registry, supplied complete source record: permanent ID \(30004342\), title ``Harder Congruences for Holomorphic Jacobi Forms,'' source problem code \texttt{catalogue entry}, in the Oberwolfach Reports Open Problems collection. The record gives proposed year 2019 and no named proposer.
\item Supplied report citation: \emph{Modular Forms (2020)}, Oberwolfach Reports, pp.\ 3529--3606, DOI \texttt{10.4171/owr/2019/57}. \url{https://doi.org/10.4171/owr/2019/57}. The background describes this report as documenting workshop discussions and open problems; no exact contribution, page, or conjecture number is supplied for this formulation.
\item Supplied literature-source attribution: Katsurada, ``holomorphic Jacobi-form Harder conjecture.'' \url{https://doi.org/10.4171/owr/2019/57}.
\item Supplied literature-source attribution: Atobe--Chida--Ibukiyama--Katsurada--Yamauchi, ``Harder's conjecture I.'' \url{https://doi.org/10.2969/jmsj/87988798}.
\item Literature assessment checked on \texttt{2026-08-22}: supplied dated curation, not independently verified. It labels the problem partially solved and reports that the holomorphic Jacobi-form version is proved in the report only under explicit arithmetic hypotheses; later work proves further cases of related Harder congruences, while the curation establishes no general unconditional resolution. The source itself requests independent verification of time-sensitive status and cited claims.
\end{itemize}
\end{document}
